% Incognita paper P006, kind `definitions`: the setting paper of the roadmap BatchedBandits.
% Rules: AGENTS.md, section 5; docs/writing.md. One source, two builds (AGENTS.md 5.2), after
% `python3 scripts/certify.py bib`:
%   reader:    latexmk -pdf -interaction=nonstopmode -halt-on-error -cd papers/P006-batched-setting/paper.tex
%   blueprint: latexmk -pdf -interaction=nonstopmode -halt-on-error -jobname=paper-blueprint
%                -usepretex='\def\IncognitaBlueprint{}' -cd papers/P006-batched-setting/paper.tex
% Planned Lean names of the nodes: DOSSIER.md (never \lean here). Never \leanok or \mathlibok.
\documentclass[11pt]{article}
\usepackage[T1]{fontenc}
\usepackage[utf8]{inputenc}
\usepackage{lmodern}
\usepackage[a4paper,margin=2.5cm]{geometry}
\usepackage{amsmath,amssymb,mathtools}
\usepackage{booktabs}
\usepackage{algorithm}
\usepackage{algpseudocode}
\usepackage{enumitem}
\usepackage{xcolor}
\usepackage[colorlinks=true,linkcolor=blue!50!black,citecolor=blue!50!black,urlcolor=blue!50!black]{hyperref}
\usepackage[round]{natbib}
\usepackage{../../templates/paper/blueprint-macros}

% Navigation: an anchor \dt{s:...} where a quantity is defined, one macro per symbol (\sy).
\newcommand{\dt}[1]{\hypertarget{#1}{}}
\newcommand{\sy}[2]{{\hyperlink{#1}{#2}}}
\newcommand{\term}[2]{\hyperlink{#1}{#2}}
\newcommand{\E}{\mathbb{E}}
\newcommand{\1}{\mathbf{1}}
\newcommand{\R}{\mathbb{R}}
\newcommand{\Nat}{\mathbb{N}}
% ---------------------------------------------------------------- arms, laws, instances
\newcommand{\Arms}{\sy{s:inst}{\mathcal A}}
\newcommand{\Kar}{\sy{s:inst}{K}}
\newcommand{\Mu}{\sy{s:inst}{\mathcal M}}
\newcommand{\mean}{\sy{s:inst}{m}}
\newcommand{\ms}{\sy{s:inst}{m^\star}}
\newcommand{\Dg}{\sy{s:inst}{\Delta}}
\newcommand{\cP}{\sy{s:inst}{\mathcal P}}
% ---------------------------------------------------------------- divergences
\newcommand{\KL}{\sy{s:kinf}{\mathrm{KL}}}
\newcommand{\Ber}{\sy{s:kinf}{\mathrm{Ber}}}
\newcommand{\kl}{\sy{s:kinf}{\mathrm{kl}}}
\newcommand{\Ki}{\sy{s:kinf}{\mathcal K_{\inf}}}
\newcommand{\Ka}{\sy{s:kinf}{K}}
% ---------------------------------------------------------------- learners, counts, regret
\newcommand{\Hs}{\sy{s:learner}{\mathcal H}}
\newcommand{\Pnu}{\sy{s:learner}{\mathbb P}}
\newcommand{\Enu}{\sy{s:learner}{\mathbb E}}
\newcommand{\Ac}{\sy{s:learner}{A}}
\newcommand{\Rw}{\sy{s:learner}{R}}
\newcommand{\Ng}{\sy{s:count}{N}}
\newcommand{\Reg}{\sy{s:count}{\mathcal R}}
% ---------------------------------------------------------------- grids
\newcommand{\Gr}{\sy{s:grid}{\mathcal G}}
\newcommand{\tl}{\sy{s:grid}{t}}
\newcommand{\Bg}{\sy{s:grid}{B}}
\newcommand{\bt}{\sy{s:grid}{b}}
\newcommand{\Lg}{\sy{s:log}{L}}
\newcommand{\xg}{\sy{s:log}{x}}
\newcommand{\Sq}{\sy{s:exgrid}{\mathsf S}}
\newcommand{\Ge}{\sy{s:exgrid}{\mathsf G}}
\newcommand{\Db}{\sy{s:exgrid}{\mathsf D}}
% ---------------------------------------------------------------- batched learners, models
\newcommand{\mask}{\sy{s:batched}{\mathrm{mask}}}
\newcommand{\Mk}[1]{\sy{s:info}{\mathcal M}_{#1}}
% ---------------------------------------------------------------- example learners
\newcommand{\RR}{\hyperref[def:exlearners]{\texttt{RR}}}
\newcommand{\Const}{\hyperref[def:exlearners]{\texttt{Const}}}
\newcommand{\FL}{\hyperref[def:exlearners]{\texttt{FL}}}
\newcommand{\React}{\hyperref[def:exlearners]{\texttt{React}}}
\newcommand{\Par}{\hyperref[def:exlearners]{\texttt{Par}}}
\newcommand{\UCB}{\hyperref[def:exlearners]{\texttt{UCB}}}
% ---------------------------------------------------------------- words
\newcommand{\sees}{\term{t:sees}{sees the rewards up to}}
\newcommand{\ignores}{\term{t:sees}{ignores the rewards}}
\newcommand{\batched}{\term{t:batched}{batched}}
\newcommand{\consistent}{\term{t:cons}{consistent at the batch ends}}
\newcommand{\seqcons}{\term{t:seqc}{sequentially consistent}}
\newcommand{\finer}{\term{t:finer}{finer}}
% ---------------------------------------------------------------- proofs
\newcommand{\Setting}{\noindent\textbf{Setting.}\ }
\newcommand{\Goal}{\par\smallskip\noindent\textbf{Goal.}\ }
\newcommand{\why}[1]{\text{\small(#1)}}
\setlength{\emergencystretch}{3em}

\title{Stochastic bandits on an imposed grid of batches:\\ the setting}
\author{Claude Opus 5.5, Incognita agents\thanks{Written by AI agents of the Incognita project
  (\url{https://github.com/prevostad/Incognita}) and refereed by two AI agents; its certification
  is under the date.}}
\date{2026-11-16}

\begin{document}
\maketitle

\begin{abstract}
We fix the definitions of a model of $K$-armed stochastic bandits in which the learner observes
the rewards only at the ends of batches, and the ends of the batches $t_1<t_2<\cdots$ are imposed
by the environment rather than chosen by the learner. The rewards are in $[0,1]$ with a largest
mean below $1$, and the relevant divergence is the non-parametric $\mathcal K_{\inf}$. We define
imposed grids, the growth ratios $x_\ell=\log t_\ell/\log t_{\ell-1}$, batched learners, a scale
of information models $\mathcal M_k$ (during batch $\ell+1$ the learner knows $t_1,\dots,t_{\ell+k}$),
and consistency at the batch ends. We prove the basic relations between these objects: the
information models form a strictly increasing chain from $\mathcal M_0$ (nothing known ahead) to
$\mathcal M_\infty$ (the whole grid known); refining a grid preserves batching; on the grid whose
batches after the first have size one, batched learners reduce to the classical sequential ones
except in the first batch, where they must ignore the rewards, and consistency at the batch ends
reduces to classical consistency; and $2\Delta_a^2\le\mathcal K_{\inf}(\nu_a,m^\star)\le\log\frac{1-m_a}{1-m^\star}$
for every suboptimal arm. Every definition comes with an explicit object that satisfies it and one
that it rules out. The paper proves no regret bound: it is the common ground on which such bounds
are stated.
\end{abstract}

% ======================================================================
\section{Introduction}\label{sec:intro}

In clinical trials, A/B tests, or any system that redeploys a policy periodically, a bandit learner
can update its decisions only at the end of batches, and when the batches end is decided by the
environment: a trial phase closes, a season ends, a new day of data arrives. The learner still
chooses an arm at every round, but it sees the rewards of a batch only once the batch is over, and
it may not even know in advance when that will be. The question that motivates this model is how
much the problem-dependent asymptotic regret of Lai, Robbins, Burnetas and Katehakis degrades as a
function of the imposed grid, and in which information model.

Before answering it, one must say precisely what a learner on an imposed grid is, what it knows of
the grid, and what consistency means when the learner can be evaluated only at the ends of the
batches. This paper fixes these definitions, proves their basic relations, and checks that none of
them is empty or trivially satisfied. It contains no regret bound.

\paragraph{The objects} An imposed grid is a strictly increasing sequence of batch ends
$0=t_0<t_1<t_2<\cdots$ with $t_1\ge\max(2,K)$; its growth is measured by
\begin{equation*}
  x_\ell := \frac{\log t_\ell}{\log t_{\ell-1}}\quad(\ell\ge2).
\end{equation*}
A learner is \emph{batched} on a grid when its action at a round of batch $\ell+1$ depends on the
rewards only through those of the rounds $1,\dots,t_\ell$. In the information model $\mathcal M_k$,
during batch $\ell+1$ the learner may also use the grid, but only through $t_1,\dots,t_{\ell+k}$:
$\mathcal M_0$ knows nothing ahead, $\mathcal M_1$ knows the end of the current batch, and
$\mathcal M_\infty$ knows the whole grid. A learner is \emph{consistent at the batch ends} on a class of
instances if $\E N_c(t_\ell)=o(t_\ell^\alpha)$ for every $\alpha>0$ and every suboptimal arm $c$.

\paragraph{Main result, informally} Theorem~\ref{thm:main} shows that this setting is coherent and
extends the classical one: the chain $\mathcal M_0\subsetneq\mathcal M_1\subsetneq\cdots\subsetneq\mathcal M_\infty$
is strictly increasing; batching is preserved when a grid is refined; and on the grid
$t_\ell=s+\ell-1$, whose batches after the first have size one, a learner is batched exactly when
its first $s$ actions ignore the rewards, and consistency at the batch ends is the classical
consistency of \citet[Definition~16.1]{LS}. Proposition~\ref{prop:kinf} brackets the divergence
that governs the problem-dependent constants: for every instance with rewards in $[0,1]$ and
largest mean $m^\star<1$, and every suboptimal arm $a$,
\begin{equation*}
  0<2\Delta_a^2\ \le\ \mathcal K_{\inf}(\nu_a,m^\star)\ \le\ \log\frac{1-m_a}{1-m^\star}<\infty .
\end{equation*}

\paragraph{Our contributions} We summarize our contributions below:
\begin{itemize}
\item precise definitions of imposed grids, batched learners, the information models
  $\mathcal M_k$ for $k\in\Nat\cup\{\infty\}$ and consistency at the batch ends, on top of the
  classical model of learners as sequences of probability kernels
  (Definitions~\ref{def:instance}--\ref{def:consistent}), with the translation to rounds numbered
  from zero fixed once (Section~\ref{sec:setting});
\item their basic relations (Theorem~\ref{thm:main}) and the bounds on $\mathcal K_{\inf}$ on
  $[0,1]$, with its value on Bernoulli laws (Proposition~\ref{prop:kinf});
\item for every definition, an explicit object that satisfies it and one that it excludes
  (Proposition~\ref{prop:witnesses} and Table~\ref{tab:witnesses}), and a learner that is consistent
  at the batch ends on the whole class of instances, on the sequential grid
  (Corollary~\ref{cor:ucb}).
\end{itemize}

\paragraph{Outline} Section~\ref{sec:related} places the model among batched and delayed bandits.
Section~\ref{sec:setting} gives the definitions. Section~\ref{sec:main} states the basic relations.
Section~\ref{sec:overview} explains how the proofs work. Section~\ref{sec:examples} gives the
witnesses and exclusions and numerical values on explicit instances, and Section~\ref{sec:open}
concludes. The proofs are in the appendices.

% ======================================================================
\section{Related work and context}\label{sec:related}

\paragraph{Sequential bandits and consistency} In the classical model the learner observes each
reward before its next decision. \citet{LR85}, for parametric families, and \citet{BK96}, for
more general models, show that a learner whose regret is $o(T^\alpha)$ for every $\alpha>0$ on
every instance of a class must pull every suboptimal arm $a$ at least $(1+o(1))\log T$ times the
inverse of a divergence. For the model of all laws on $[0,1]$ this divergence is
$\mathcal K_{\inf}(\nu_a,m^\star)$, and the bound is \citet[Theorem~1]{GMS19}. We use the definition
of consistency of \citet[Definition~16.1]{LS} (regret $o(T^\alpha)$) and, equivalently for finitely
many arms, the ``uniformly fast convergent'' strategies of \citet[Definition~1]{GMS19} (counts
$o(T^\alpha)$). We model a learner as a sequence of probability kernels, which carry its internal
randomisation; \citet[Section~1.1]{GMS19} make that randomisation explicit with independent
uniform variables instead.

\paragraph{The divergence $\mathcal K_{\inf}$ on $[0,1]$} For distributions on $[0,1]$ the
relevant divergence is the infimum of $\mathrm{KL}(\nu_a,\nu')$ over the laws $\nu'$ on $[0,1]$ with
mean above $m^\star$ \citep{BK96}, \citep[Section~1.2]{GMS19}. \citet[Theorem~8]{HT10} prove a dual
formula for it, which algorithms such as IMED use (see also \citet{HT15}). Empirical KL-UCB attains
the constant $1/\mathcal K_{\inf}$ for finitely supported laws on $[0,1]$, when every mean is
positive and the largest one is below $1$ \citep[Theorem~2]{CGMMS13}. \citet[Section~6.1]{CGMMS13}
also observe that $\mathcal K_{\inf}$ is at least the Bernoulli divergence of the means, the lower
bound of our Proposition~\ref{prop:kinf}. The dual formula is not needed for the definitions; we
prove here only what follows from the primal definition.

\paragraph{Batched bandits with a chosen grid} \citet{PRCS16} study policies that must split the
trials into a small number of batches whose ends the policy designs, and give close to minimax
optimal regret bounds; \citet{GHRZ19} extend this to $K$ arms, with lower bounds that hold even when
the batch sizes are chosen adaptively. \citet[Theorem~4]{JTXHXG21} consider anytime learners on a
static grid and show that, for one-parameter exponential families, no anytime algorithm reaches the
sequential constant with $c\log\log T$ batches. See also \citet{EKMM21} and \citet{JXXG21}. In all
these works the grid is a design variable of the learner.

\paragraph{Imposed batches and delays} \citet{KO21b} analyse Thompson sampling when the batch sizes
are not chosen by the learner; in their model \citep[Section~II.A]{KO21b} the end of a batch may be
fixed at the previous end, possibly by an adversary, as a function of the history.
\citet{PGPK21} take batches of a common size $b$ fixed by the environment and, under two
assumptions on the policy, bound the regret of its batched version over $n=bM$ rounds between the
regret of the policy run online and $b$ times its regret over $M$ rounds
\citep[Theorem~3.1]{PGPK21}; see also \citet{PGPK22} for linear bandits.
\citet{QW26} let the environment fix the update times $t_1<\dots<t_B=T$
adversarially; a \emph{batch-aware} learner is given all of them in advance, and a
\emph{batch-oblivious} learner is only told when each batch ends. For stochastic bandits with
rewards in $[0,1]$ they measure the hardness of the grid by $\alpha=\max_j(t_j-t_{j-1})/t_{j-1}$
(with $t_0:=1$), prove a regret bound of order $\alpha^2\log(KT)\sum_a1/\Delta_a$ and a lower bound
of order $\alpha^{3/2}\sum_a1/\Delta_a$ on some instances \citep[Theorems~1 and~4]{QW26}. Batching is
also a form of delayed feedback: the reduction of \citet[Theorem~6]{JGS13} adds to the regret
without delays a term proportional to the largest number of missing observations, here the size
of the largest batch.

\paragraph{Closest related work} Table~\ref{tab:comparison} compares the batch models. The closest
is that of \citet{QW26}: the grid is imposed by the environment, and their two kinds of learners
are the two ends of our scale, the batch-aware learner corresponding to $\mathcal M_\infty$ (the
whole grid known) and the batch-oblivious one to $\mathcal M_0$ (only the ends already passed
known), up to the horizon $T$, which their algorithm takes as an input. It does not give our results: it has no intermediate model between the two, its criterion
is the regret at a finite horizon $T$, up to constant factors, with no consistency notion and no
problem-dependent constant such as $1/\mathcal K_{\inf}$, and its hardness $\alpha$ is a ratio of
consecutive ends, while the asymptotic constants of the classical theory involve ratios of their
logarithms, our $x_\ell$. In \citet{KO21b} the batch ends are imposed too, but the end of a batch is
decided at the previous end and may depend on the history, and the algorithm analysed, Thompson
sampling, does not use the batch ends; the criterion is the regret over all rounds. In
\citet{PGPK21} the batch size is imposed and constant. The chosen-grid models of \citet{PRCS16},
\citet{GHRZ19} and \citet{JTXHXG21} do not contain imposed grids: there the learner picks the ends,
which gives it at least the information of $\mathcal M_\infty$ together with the freedom to choose
the grid.

\paragraph{Gap with the literature} Imposed grids, and learners that know the whole grid or only
its past ends, are in \citet{QW26}. We are not aware of a published definition of the intermediate
information models $\mathcal M_k$, $1\le k<\infty$, of consistency at the batch ends, or of a
statement relating them to the sequential notions. This paper provides them, for rewards in
$[0,1]$, together with witnesses and exclusions.

\begin{table}[t]
\centering
\caption{Batch models: who fixes the batch ends, and what the learner knows of them.}\label{tab:comparison}
\begin{tabular}{@{}p{0.2\textwidth}p{0.22\textwidth}p{0.3\textwidth}p{0.2\textwidth}@{}}
\toprule
Reference & Batch ends fixed by & Known to the learner & Results \\
\midrule
\citet{PRCS16} & the learner & all (designed) & minimax regret \\
\citet{GHRZ19} & the learner, possibly adaptively & all (designed) & minimax regret \\
\citet{JTXHXG21} & the learner (static grid) & all; the horizon unknown & asymptotic regret \\
\citet{KO21b} & the environment, at the previous end & not used by the algorithm & asymptotic regret of Thompson sampling \\
\citet{PGPK21} & the environment (constant size $b$) & not needed: the policy updates when the rewards arrive & regret relative to online regret \\
\citet{QW26} & the environment, adversarially & all (batch-aware) or the past ends (batch-oblivious) & regret at horizon $T$, up to constants \\
This paper & the environment, in advance & $t_1,\dots,t_{\ell+k}$ during batch $\ell+1$ (model $\mathcal M_k$) & definitions only \\
\bottomrule
\end{tabular}
\end{table}

% ======================================================================
\section{Problem formulation}\label{sec:setting}

We first recall the classical sequential setting (instances, the divergence $\Ki$, learners,
counts, regret, consistency), then define imposed grids, batched learners, the information models
and consistency at the batch ends. Logarithms are natural.

\subsection{Classical setting}

\begin{definition}[{laws on $[0,1]$, instances, the class $\mathcal P$}]\label{def:instance}
\dt{s:inst}Let $\Mu$ be the set of probability measures $Q$ on $\R$ (with its Borel
$\sigma$-algebra) such that $Q([0,1])=1$, and for $Q\in\Mu$ let $\mean(Q):=\int x\,Q(dx)\in[0,1]$
be its mean. The arms form a finite set $\Arms$ with $\Kar:=|\Arms|\ge2$ elements. An instance is
a family $\nu=(\nu_c)_{c\in\Arms}$ of probability measures on $\R$; when every $\nu_c$ is in $\Mu$, we
write
\begin{equation*}
  \mean_c:=\mean(\nu_c),\qquad \ms:=\max_{c\in\Arms}\mean_c,\qquad \Dg_c:=\ms-\mean_c\ge0 .
\end{equation*}
An arm $c$ is \emph{optimal} if $\Dg_c=0$ and \emph{suboptimal} if $\Dg_c>0$. The class
$\cP:=\{\nu:\ \nu_c\in\Mu\ \text{for all}\ c\in\Arms,\ \ms<1\}$ is the set of instances with
rewards in $[0,1]$ and largest mean below $1$.
\end{definition}

\begin{definition}[divergences and $\mathcal K_{\inf}$]\label{def:kinf}
\uses{def:instance}
\dt{s:kinf}For probability measures $P,Q$ on a measurable space,
$\KL(P,Q):=\int\log\frac{dP}{dQ}\,dP\in[0,\infty]$ if $P$ is absolutely continuous with respect to
$Q$, and $\KL(P,Q):=+\infty$ otherwise. For $p\in[0,1]$, $\Ber(p):=(1-p)\delta_0+p\,\delta_1\in\Mu$,
and $\kl(p,q):=\KL(\Ber(p),\Ber(q))$, that is
\begin{equation*}
  \kl(p,q)=p\log\frac pq+(1-p)\log\frac{1-p}{1-q}\qquad(p,q\in[0,1]),
\end{equation*}
with $0\log\frac0q:=0$ for every $q\in[0,1]$ (so $\kl(p,p)=0$) and $\kl(p,q):=+\infty$ when
$q\in\{0,1\}$ and $q\ne p$. For $P\in\Mu$ and a level $\mu\in\R$,
\begin{equation*}
  \Ki(P,\mu):=\inf\big\{\KL(P,Q):\ Q\in\Mu,\ \mean(Q)>\mu\big\}\in[0,\infty],
\end{equation*}
with $\inf\emptyset:=+\infty$. For an instance $\nu\in\cP$ and an arm $a$ we write
$\Ka_a:=\Ki(\nu_a,\ms)$.
\end{definition}

The infimum is over all laws on $[0,1]$, not over a parametric family. On Bernoulli laws it
coincides with $\kl$ (Proposition~\ref{prop:kinf}(iv)).

\begin{definition}[learners and the law of the interaction]\label{def:learner}
\uses{def:instance}
\dt{s:learner}Rounds are numbered $t=1,2,\dots$ For $t\ge0$, a history of length $t$ is an element
$h=(a_1,y_1,\dots,a_t,y_t)$ of $\Hs_t:=(\Arms\times\R)^t$ ($\Hs_0$ is a point), with the product
$\sigma$-algebra ($\Arms$ discrete). A \emph{learner} is a sequence $\pi=(\pi_t)_{t\ge1}$ where
$\pi_t$ is a probability kernel from $\Hs_{t-1}$ to $\Arms$: $\pi_t(\cdot\mid h)$ is the law of the
arm pulled at round $t$ after the history $h$. It is \emph{deterministic} if every
$\pi_t(\cdot\mid h)$ is a Dirac mass. For an instance $\nu$ and a learner $\pi$, $\Pnu_{\nu,\pi}$ is
the unique probability measure on $(\Arms\times\R)^{\Nat}$, with coordinates
$(\Ac_t,\Rw_t)_{t\ge1}$, under which, for every $t\ge1$,
\begin{equation*}
  \Ac_t\ \big|\ H_{t-1}\ \sim\ \pi_t(\cdot\mid H_{t-1}),\qquad
  \Rw_t\ \big|\ (H_{t-1},\Ac_t)\ \sim\ \nu_{\Ac_t},
\end{equation*}
where $H_{t}:=(\Ac_1,\Rw_1,\dots,\Ac_t,\Rw_t)$; it exists and is unique by the Ionescu-Tulcea
theorem. $\Enu_{\nu,\pi}$ is its expectation.
\end{definition}

The internal randomisation of the learner is carried by the kernels $\pi_t$. Without further
constraint, the policy of round $t$ may use every reward of the rounds $1,\dots,t-1$: this is the
classical sequential setting, and a learner in this sense is also called \emph{sequential}.

\begin{definition}[counts and regret]\label{def:counts}
\uses{def:instance, def:learner}
\dt{s:count}For an arm $c$ and $t\ge0$, the number of pulls of $c$ up to round $t$ included is
\begin{equation*}
  \Ng_c(t):=\sum_{s=1}^{t}\1(\Ac_s=c),
\end{equation*}
and, for $\nu$ with every $\nu_c\in\Mu$, the regret at round $t$ is
$\Reg_{\nu,\pi}(t):=\sum_{c\in\Arms}\Dg_c\,\Enu_{\nu,\pi}[\Ng_c(t)]$.
\end{definition}

\begin{lemma}[counts and regret]\label{lem:counts}
\uses{def:instance, def:learner, def:counts}
For every learner $\pi$, every instance $\nu$ with $\nu_c\in\Mu$ for all $c$, and every $t\ge0$:
\begin{enumerate}[label=(\roman*)]
\item $0\le\Ng_c(t)\le t$, $t\mapsto\Ng_c(t)$ is nondecreasing, and $\sum_{c\in\Arms}\Ng_c(t)=t$;
\item $\Reg_{\nu,\pi}(t)=\Enu_{\nu,\pi}\big[\sum_{s=1}^{t}\Dg_{\Ac_s}\big]=t\,\ms-\Enu_{\nu,\pi}\big[\sum_{s=1}^t\mean_{\Ac_s}\big]$,
  and $0\le\Reg_{\nu,\pi}(t)\le t\max_c\Dg_c\le t$.
\end{enumerate}
\end{lemma}
(see the proof in Appendix~\ref{sec:pf-lem:counts}).

\begin{definition}[sequential consistency]\label{def:seqcons}
\uses{def:instance, def:learner, def:counts}
\dt{t:seqc}Let $E\subseteq\cP$. A learner $\pi$ is \emph{sequentially consistent} on $E$ if, for
every $\nu\in E$ and every $\alpha>0$,
\begin{equation*}
  \lim_{T\to\infty}\frac{\Reg_{\nu,\pi}(T)}{T^{\alpha}}=0 .
\end{equation*}
\end{definition}

This is \citet[Definition~16.1]{LS}, restricted to a class of instances with rewards in $[0,1]$.

\subsection{Imposed grids}

\begin{definition}[grid, batches]\label{def:grid}
\uses{def:instance}
\dt{s:grid}A \emph{grid} (for $\Kar$ arms) is a sequence $g=(\tl_\ell)_{\ell\ge0}$ of integers with
\begin{equation*}
  \tl_0=0,\qquad \tl_{\ell-1}<\tl_\ell\ \ (\ell\ge1),\qquad \tl_1\ge\max(2,\Kar).
\end{equation*}
$\Gr$ is the set of grids. Batch $\ell\ge1$ is the set of rounds $\{\tl_{\ell-1}+1,\dots,\tl_\ell\}$,
its size is $\Bg_\ell:=\tl_\ell-\tl_{\ell-1}\ge1$, and $\tl_\ell$ is its \emph{end}. For a round
$t\ge1$, its batch is $\bt(t):=\min\{\ell\ge1:\ t\le\tl_\ell\}$; when several grids are involved,
$\bt_g(t)$ denotes the batch of $t$ in the grid $g$. A grid $g$ is \dt{t:finer}\emph{finer}
than a grid $g'$ if every end of $g'$ is an end of $g$.
\end{definition}

The condition $\tl_1\ge2$ makes $\log\tl_1>0$, so that the ratios below are defined; the condition
$\tl_1\ge\Kar$ lets the first batch pull every arm once.

\begin{definition}[logarithmic ends and ratios]\label{def:logratio}
\uses{def:grid}
\dt{s:log}For a grid $g=(\tl_\ell)_\ell$, $\Lg_\ell:=\log\tl_\ell$ for $\ell\ge1$, $\xg_1:=1$ and
\begin{equation*}
  \xg_\ell:=\frac{\Lg_\ell}{\Lg_{\ell-1}}\qquad(\ell\ge2).
\end{equation*}
\end{definition}

\begin{lemma}[elementary facts on grids]\label{lem:grid}
\uses{def:grid, def:logratio}
For every grid $g$:
\begin{enumerate}[label=(\roman*)]
\item $\tl_\ell\ge\tl_1+\ell-1$ for every $\ell\ge1$; in particular $\tl_\ell\to\infty$;
\item every round $t\ge1$ belongs to exactly one batch, namely $\bt(t)$, and
  $\tl_{\bt(t)-1}<t\le\tl_{\bt(t)}$; moreover $\bt(\tl_\ell)=\ell$, $\bt$ is nondecreasing, and
  $\bt(t)=1$ if and only if $t\le\tl_1$;
\item $\Lg_\ell\ge\log2>0$ for $\ell\ge1$, $\ell\mapsto\Lg_\ell$ is strictly increasing,
  $\xg_\ell>1$ for $\ell\ge2$, and for $\ell\ge2$
  \begin{equation*}
    \Lg_\ell=\Lg_1\prod_{k=2}^{\ell}\xg_k,\qquad \tl_\ell=\tl_{\ell-1}^{\,\xg_\ell}.
  \end{equation*}
\end{enumerate}
\end{lemma}
(see the proof in Appendix~\ref{sec:pf-lem:grid}).

\begin{definition}[three example grids]\label{def:exgrids}
\uses{def:grid}
\dt{s:exgrid}For an integer $s\ge\max(2,\Kar)$ and $\ell\ge1$, with $\tl_0:=0$ in each case:
\begin{equation*}
  \Sq_s:\ \tl_\ell:=s+\ell-1,\qquad
  \Ge_s:\ \tl_\ell:=s\,2^{\ell-1},\qquad
  \Db_s:\ \tl_\ell:=s^{2^{\ell-1}} .
\end{equation*}
$\Sq_s$ is the \emph{sequential grid}: its batches after the first have size one.
\end{definition}

\paragraph{Rounds numbered from zero} A formalisation of these definitions in the proof assistant
Lean numbers the rounds from $0$, as the Lean library of learning algorithms on which it is built
does: the round $t\ge1$ of this paper is the round $t-1$ there. With this shift, batch $\ell$ is the set of rounds $\tl_{\ell-1},\dots,\tl_\ell-1$,
the batch of a round $n\ge0$ is $\min\{\ell\ge1:\ n<\tl_\ell\}$, the count $\Ng_c(t)$ is the
number of pulls of $c$ in the rounds $0,\dots,t-1$, and the rewards available during batch $\ell+1$
are those of the rounds $<\tl_\ell$. In particular every quantity evaluated at a batch end
$\tl_\ell$ (counts, regret) is the same in both numberings, and no statement at the batch ends
needs a shift.

\subsection{Batched learners and information models}

\begin{definition}[batched learner]\label{def:batched}
\uses{def:learner, def:grid}
\dt{s:batched}Let $t\ge1$ and let $r\ge0$ be an integer. For
$h=(a_1,y_1,\dots,a_{t-1},y_{t-1})\in\Hs_{t-1}$, let
$\mask_r(h):=(a_1,y_1\1(1\le r),\dots,a_{t-1},y_{t-1}\1(t-1\le r))$: the rewards of the rounds
$s>r$ are replaced by $0$. The policy $\pi_t$ \dt{t:sees}\emph{sees the rewards up to} $r$ if
\begin{equation*}
  \pi_t(\cdot\mid h)=\pi_t(\cdot\mid h')
\end{equation*}
for all $h,h'\in\Hs_{t-1}$ with $a_s=a'_s$ for every $s<t$ and $y_s=y'_s$ for every $s\le\min(r,t-1)$;
it \emph{ignores the rewards} if it \sees{} $0$. A learner $\pi$ is \dt{t:batched}\emph{batched} on
a grid $g$ if, for every round $t\ge1$,
\begin{equation*}
  \pi_t\ \text{\sees}\ \tl_{\bt(t)-1}.
\end{equation*}
\end{definition}

So during batch $\ell+1$ the learner uses its own past actions, its internal randomisation and the
rewards of the rounds $1,\dots,\tl_\ell$, and nothing of the rewards of the current batch. A learner
batched on $g$ may depend on $g$ in any way: on a fixed grid, it knows the whole grid.

\begin{lemma}[masked form]\label{lem:mask}
\uses{def:batched, def:learner, def:grid}
Let $\pi$ be a learner. For $t\ge1$ and an integer $r\ge0$, $\mask_r$ is measurable, and $\pi_t$
\sees{} $r$ if and only if $\pi_t(\cdot\mid h)=\pi_t(\cdot\mid\mask_r(h))$ for every
$h\in\Hs_{t-1}$. Hence, for a grid $g=(\tl_\ell)_\ell$, $\pi$ is
\batched{} on $g$ if and only if $\pi_t$ is the kernel $h\mapsto\pi_t(\cdot\mid\mask_{\tl_{\bt(t)-1}}(h))$
for every $t\ge1$.
\end{lemma}
(see the proof in Appendix~\ref{sec:pf-lem:mask}).

\begin{definition}[information models]\label{def:infomodel}
\uses{def:grid, def:learner, def:batched}
\dt{s:info}A \emph{learner for imposed grids} is a family $\Pi=(\pi^g)_{g\in\Gr}$ of learners, one
per grid. For $k\in\Nat\cup\{\infty\}$, $\Pi$ belongs to the information model $\Mk{k}$ if
\begin{enumerate}[label=(\alph*)]
\item for every grid $g$, $\pi^g$ is \batched{} on $g$;
\item for all grids $g=(\tl_j)_j$ and $g'=(\tl'_j)_j$, every $\ell\ge0$ and every round $t$ that
  belongs to batch $\ell+1$ of $g$ and to batch $\ell+1$ of $g'$,
  \begin{equation*}
    \tl_j=\tl'_j\ \text{for every integer}\ j\ \text{with}\ 1\le j\le\ell+k
    \quad\Longrightarrow\quad\pi^g_t=\pi^{g'}_t .
  \end{equation*}
\end{enumerate}
\end{definition}

During batch $\ell+1$, a learner of $\Mk{k}$ uses the grid only through $\tl_1,\dots,\tl_{\ell+k}$,
the current round and the fact that the batch has not ended. In $\Mk{0}$ it knows nothing ahead
(neither the end of the current batch nor any horizon: the doubly anytime model); in $\Mk{1}$ it
knows the end of the current batch; for $k=\infty$, the hypothesis of (b) forces $g=g'$, so (b) is
void and $\Mk{\infty}$ is the set of families of batched learners, which may know the whole grid.

\subsection{Consistency at the batch ends}

\begin{definition}[consistency at the batch ends]\label{def:consistent}
\uses{def:instance, def:grid, def:learner, def:counts}
\dt{t:cons}Let $E\subseteq\cP$ and let $g=(\tl_\ell)_\ell$ be a grid. A learner $\pi$ is \emph{consistent at the
batch ends} of $g$ on $E$ if, for every $\nu\in E$, every suboptimal arm $c$ of $\nu$ and every
$\alpha>0$,
\begin{equation*}
  \lim_{\ell\to\infty}\frac{\Enu_{\nu,\pi}[\Ng_c(\tl_\ell)]}{\tl_\ell^{\,\alpha}}=0 .
\end{equation*}
\end{definition}

The criterion is evaluated only at the ends of the batches, where the learner has seen every
reward, which is natural when it cannot react inside a batch. Written with counts, as here, it is
the ``uniformly fast convergent'' property of \citet[Definition~1]{GMS19} along the batch ends;
written with the regret, it is the consistency of \citet[Definition~16.1]{LS} along the batch ends.
The two forms agree:

\begin{lemma}[counts or regret]\label{lem:consist-regret}
\uses{def:consistent, def:counts, def:instance, def:grid, def:learner}
Let $E\subseteq\cP$, $g=(\tl_\ell)_\ell$ a grid and $\pi$ a learner. Then $\pi$ is \consistent{} of $g$ on $E$
if and only if $\Reg_{\nu,\pi}(\tl_\ell)/\tl_\ell^{\,\alpha}\to0$ as $\ell\to\infty$ for every
$\nu\in E$ and every $\alpha>0$.
\end{lemma}
(see the proof in Appendix~\ref{sec:pf-lem:consist-regret}).

% ======================================================================
\section{Basic relations}\label{sec:main}

The classical facts are the sequential ones: a sequential learner is any learner of
Definition~\ref{def:learner}, and it is sequentially consistent in the sense of
\citet[Definition~16.1]{LS}. The first result places the batched objects with respect to them and
to each other.

\begin{theorem}[basic relations]\label{thm:main}
\uses{def:instance, def:learner, def:grid, def:exgrids, def:batched, def:infomodel, def:consistent, def:seqcons}
Fix the set of arms $\Arms$, with $\Kar\ge2$. Then:
\begin{enumerate}[label=(\roman*)]
\item \emph{(Information models.)} For all $k\le k'$ in $\Nat\cup\{\infty\}$,
  $\Mk{k}\subseteq\Mk{k'}$, and the inclusion is strict when $k<k'$. A learner $\pi$ is \batched{}
  on a grid $g$ if and only if $\pi=\pi^g$ for some $\Pi\in\Mk{\infty}$.
\item \emph{(Refinement.)} Let $g$ be a grid \finer{} than a grid $g'$ and $E\subseteq\cP$. Every
  learner \batched{} on $g'$ is \batched{} on $g$, and every learner \consistent{} of $g$ on $E$ is
  \consistent{} of $g'$ on $E$.
\item \emph{(Sequential grid.)} Let $s\ge\max(2,\Kar)$ and $E\subseteq\cP$. A learner $\pi$ is
  \batched{} on $\Sq_s$ if and only if $\pi_t$ \ignores{} for every $t\le s$; and $\pi$ is
  \consistent{} of $\Sq_s$ on $E$ if and only if it is \seqcons{} on $E$.
\end{enumerate}
\end{theorem}
(see the proof in Appendix~\ref{sec:pf-thm:main}).

\paragraph{Interpretation} Item (i) says that the information models are genuinely different: a
learner that reads $\tl_{\ell+k+1}$ during batch $\ell+1$ belongs to $\Mk{k+1}$ and not to $\Mk{k}$
(Proposition~\ref{prop:witnesses}(d)). A lower bound proved for every batched learner on a fixed
grid is therefore a lower bound in $\Mk{\infty}$, hence in every $\Mk{k}$; an upper bound for a learner
of $\Mk{0}$ holds in every $\Mk{k}$. Item (ii) says that observing more often can only help a learner:
a batched learner of a coarse grid is still batched on every finer grid. Item (iii) is the bridge to
the classical setting: on $\Sq_s$ the only constraint is that the first $s\ge\Kar$ rounds do not
react to rewards, which classical index policies satisfy when they start by pulling each arm once
and $s=\Kar$ (Corollary~\ref{cor:ucb} for UCB1); and consistency at the batch ends is then the classical notion of consistency used
by the lower bounds of \citet{LR85}, \citet{BK96} and \citet[Theorem~16.2]{LS}, which therefore
apply on $\Sq_s$ whenever the class of instances contains the alternatives their proofs use.

\paragraph{Scope} Theorem~\ref{thm:main} is about the definitions only. It does not say on which
grids a consistent learner exists, nor give any bound on the counts; it does not cover random grids
or grids that depend on the history, nor information models finer than $\Mk{k}$ (for instance a
learner that knows the sizes of the batches but not their ends).

The second result concerns the divergence that appears in the problem-dependent constants.

\begin{proposition}[{$\mathcal K_{\inf}$ on $[0,1]$}]\label{prop:kinf}
\uses{def:instance, def:kinf}
Let $P\in\Mu$ with mean $p:=\mean(P)$, and let $\mu\in\R$. Then:
\begin{enumerate}[label=(\roman*)]
\item if $\mu\ge1$, then $\Ki(P,\mu)=+\infty$;
\item if $\mu<p$, then $\Ki(P,\mu)=0$;
\item if $p\le \mu<1$, then
  \begin{equation*}
    2(\mu-p)^2\ \le\ \kl(p,\mu)\ \le\ \Ki(P,\mu)\ \le\ \log\frac{1-p}{1-\mu}<\infty ;
  \end{equation*}
\item if $P=\Ber(p)$ and $p\le \mu<1$, then $\Ki(P,\mu)=\kl(p,\mu)$;
\item for every $\nu\in\cP$ and every suboptimal arm $a$,
  $0<2\Dg_a^2\le \Ka_a\le\log\frac{1-\mean_a}{1-\ms}<\infty$.
\end{enumerate}
\end{proposition}
(see the proof in Appendix~\ref{sec:pf-prop:kinf}).

\paragraph{Why the largest mean is below one} If the largest mean is $1$, then $\Ka_a=+\infty$ for every suboptimal arm by
(i), and the classical constant $1/\Ka_a$ vanishes: such instances are excluded from $\cP$, as in
\citet[Theorem~8]{HT10}, whose dual formula holds for $\mu<1$. Item (v) shows that on $\cP$ the
constant $1/\Ka_a$ is finite and positive, so that a bound of the form $\E\Ng_a(\tl_\ell)\approx
C\log\tl_\ell/\Ka_a$ is never trivial.

% ======================================================================
\section{How the proofs work}\label{sec:overview}

All proofs are short and self-contained, except for two cited results.

\paragraph{Information models} The inclusions of Theorem~\ref{thm:main}(i) hold because agreeing
on more batch ends implies agreeing on fewer. Strictness needs, for each $k$, a learner of
$\Mk{k+1}$ that is not in $\Mk{k}$: the learner $\Par^{(k)}$, which plays one of two arms according
to the parity of $\tl_{\ell+k+1}$, separates two grids that agree up to $\tl_{\ell+k}$.

\paragraph{Refinement and the sequential grid} If every end of $g'$ is an end of $g$, the last end of
$g$ before a round $t$ is at least the last end of $g'$ before $t$; a policy that sees the rewards
up to the smaller one sees them up to the larger one. On $\Sq_s$ the last end before a round $t>s$
is $t-1$, so batching constrains only the first batch, and the batch ends $\{s,s+1,\dots\}$ are a
tail of all rounds, so limits along them are limits along all rounds; Lemma~\ref{lem:consist-regret}
passes from counts to regret.

\paragraph{The divergence} The lower bound of Proposition~\ref{prop:kinf}(iii) is the inequality
$\KL(P,Q)\ge\kl(\mean(P),\mean(Q))$ for laws on $[0,1]$, already observed by
\citet[Section~6.1]{CGMMS13}; we derive it from the data-processing inequality of
\citet[Lemma~1]{GMS19} applied to the reward itself, then use the monotonicity of $\kl$ and
Pinsker's inequality for Bernoulli laws (Lemma~\ref{lem:kl}), as in \citet[(7)]{GHMS22}. The upper
bound uses the mixture $Q_u:=(1-u)P+u\delta_1$, whose density bound gives
$\KL(P,Q_u)\le\log\frac1{1-u}$.

\paragraph{A consistent learner} Corollary~\ref{cor:ucb} combines Theorem~\ref{thm:main}(iii) with
the logarithmic regret bound of UCB1 \citep[Theorem~1]{ACF02}: UCB1 starts by pulling each arm once,
so it is batched on the sequential grid $\Sq_{\Kar}$, and logarithmic regret is sequential
consistency.

\paragraph{Map of the appendices} Appendix~\ref{sec:cited} (the cited results),
Appendix~\ref{sec:pf-setting} (the lemmas of Section~\ref{sec:setting}),
Appendices~\ref{sec:pf-thm:main}, \ref{sec:pf-prop:kinf} and~\ref{sec:pf-prop:witnesses} (the three
main statements, and Corollary~\ref{cor:ucb} in the last one).

% ======================================================================
\section{Witnesses, exclusions and examples}\label{sec:examples}

Each definition of Section~\ref{sec:setting} is checked against two explicit objects: one that
satisfies it and is not degenerate, and one that it must exclude. The learners used are the
following.

\begin{definition}[example learners]\label{def:exlearners}
\uses{def:instance, def:learner, def:grid}
\dt{s:exlearn}Enumerate the arms as $\Arms=\{c_1,\dots,c_{\Kar}\}$. All the learners below are
deterministic; we give the arm they pull at round $t$ after a history
$h=(a_1,y_1,\dots,a_{t-1},y_{t-1})$.
\begin{itemize}
\item $\RR$ (round robin) pulls $c_i$ with $i:=((t-1)\bmod\Kar)+1$.
\item $\Const_c$, for an arm $c$, always pulls $c$.
\item $\React$ pulls $c_1$ at $t=1$; at $t=2$ it pulls $c_1$ if $y_1\ge1/2$ and $c_2$ otherwise;
  for $t\ge3$ it pulls $c_1$.
\item $\FL^g$ (frozen leader), for a grid $g=(\tl_j)_j$, plays as $\RR$ if $t\le\tl_1$; otherwise,
  with $\ell:=\bt(t)-1\ge1$ and, for each arm $c$, $n_c:=|\{s\le\tl_\ell:\ a_s=c\}|$, it pulls the
  first arm (in the order $c_1,\dots,c_{\Kar}$) that maximises the frozen empirical mean
  \begin{equation*}
    \hat\mu_c:=\frac1{\max(n_c,1)}\sum_{s\le\tl_\ell:\ a_s=c}y_s .
  \end{equation*}
\item $\Par^{(k),g}$, for $k\in\Nat$ and a grid $g=(\tl_j)_j$, pulls $c_1$ if $\tl_{\bt(t)+k}$ is
  even and $c_2$ otherwise.
\item $\UCB$ (the policy UCB1 of \citet{ACF02}, with a fixed order and a fixed tie rule) pulls $c_t$
  if $t\le\Kar$; for $t>\Kar$, with, for each arm $c$, $n_c:=|\{s\le t-1:\ a_s=c\}|$ and
  $\bar y_c:=\frac1{\max(n_c,1)}\sum_{s\le t-1:\ a_s=c}y_s$, it pulls the first arm (in the order
  $c_1,\dots,c_{\Kar}$) that maximises
  \begin{equation*}
    \bar y_c+\sqrt{\frac{2\log(t-1)}{\max(n_c,1)}} .
  \end{equation*}
\end{itemize}
$\FL$ and $\Par^{(k)}$ denote the learners for imposed grids $(\FL^g)_{g\in\Gr}$ and
$(\Par^{(k),g})_{g\in\Gr}$. For an arm $c$, $\cP_c:=\{\nu\in\cP:\ \mean_c=\ms\}$.
\end{definition}

The decision rules are measurable functions of the history (finitely many comparisons of
measurable functions), so these are learners in the sense of Definition~\ref{def:learner}.

\begin{proposition}[witnesses and exclusions]\label{prop:witnesses}
\uses{def:instance, def:grid, def:logratio, def:exgrids, def:learner, def:batched, def:infomodel, def:consistent, def:exlearners}
Let $\Kar\ge2$ and $s\ge\max(2,\Kar)$.
\begin{enumerate}[label=(\alph*)]
\item \emph{Class $\cP$.} Every Bernoulli instance $(\Ber(p_c))_{c\in\Arms}$ with all $p_c<1$ is in
  $\cP$. An instance with an arm $\delta_1$, or with an arm of law $\mathcal N(0,1)$, is not in $\cP$.
\item \emph{Grids.} $\Sq_s$, $\Ge_s$ and $\Db_s$ are grids; $\xg_\ell\to1$ on $\Sq_s$ and on $\Ge_s$,
  and $\xg_\ell=2$ for every $\ell\ge2$ on $\Db_s$. The sequence $\tl_\ell:=\ell$ is not a grid.
\item \emph{Batched learners.} $\RR$ and $\Const_c$ are \batched{} on every grid. For every grid
  $g$, $\FL^g$ is \batched{} on $g$ and $\pi:=\FL^g$ does not \term{t:sees}{ignore the rewards} at
  round $\tl_1+1$. $\React$ is \batched{} on no grid.
\item \emph{Information models.} $\FL\in\Mk{0}$, and, for every arm $c$, the constant family
  $(\Const_c)_{g\in\Gr}$ is in $\Mk{0}$. For every $k\in\Nat$, $\Par^{(k)}\in\Mk{k+1}$ and
  $\Par^{(k)}\notin\Mk{k}$.
\item \emph{Consistency.} For every arm $c$ and every grid, $\Const_c$ is \consistent{} on $\cP_c$.
  On every grid, neither $\Const_c$ nor $\RR$ is \consistent{} on $\cP$.
\end{enumerate}
\end{proposition}
(see the proof in Appendix~\ref{sec:pf-prop:witnesses}).

The witness of consistency in (e) is degenerate: on $\cP_c$ nothing has to be learned. A learner
consistent at the batch ends on the whole class $\cP$ follows from Theorem~\ref{thm:main}(iii) and
the classical regret bound of UCB1 \citep[Theorem~1]{ACF02} (Theorem~\ref{cite:ACF02-thm1}).

\begin{corollary}[a consistent learner on the sequential grid]\label{cor:ucb}
\uses{def:instance, def:exgrids, def:batched, def:consistent, def:exlearners}
Let $\Kar\ge2$. The learner $\UCB$ is \batched{} on $\Sq_{\Kar}$ and \consistent{} of $\Sq_{\Kar}$
on $\cP$.
\end{corollary}
(see the proof in Appendix~\ref{sec:pf-cor:ucb}).

On a grid with batches of size larger than one after the first, whether a learner consistent at the
batch ends on $\cP$ exists, and in which information model, is not settled here.

Together with Proposition~\ref{prop:kinf}(i) and~(iv) for $\Ki$, this gives a witness and an
exclusion for each definition (Table~\ref{tab:witnesses}). The exclusions are the ones the
definitions are meant to make: a learner that reacts inside a batch, a learner that reads the future
of the grid, a learner that never explores, an instance whose best mean is $1$.

\begin{table}[t]
\centering
\caption{A witness and an exclusion for each definition.}\label{tab:witnesses}
\begin{tabular}{@{}p{0.22\textwidth}p{0.36\textwidth}p{0.34\textwidth}@{}}
\toprule
Definition & Witness & Exclusion \\
\midrule
class $\cP$ (Def.~\ref{def:instance}) & Bernoulli instances with means $<1$ & an arm $\delta_1$ ($\ms=1$); an arm $\mathcal N(0,1)$ \\
$\Ki$ (Def.~\ref{def:kinf}) & $\Ki(\Ber(p),\mu)=\kl(p,\mu)$, $p\le\mu<1$ & $\Ki(P,\mu)=+\infty$ for $\mu\ge1$ \\
grid (Def.~\ref{def:grid}) & $\Sq_s$, $\Ge_s$ ($\xg_\ell\to1$), $\Db_s$ ($\xg_\ell=2$) & $\tl_\ell=\ell$ \\
batched (Def.~\ref{def:batched}) & $\FL^g$ on $g$ (uses rewards) & $\React$, on every grid \\
$\Mk{k}$ (Def.~\ref{def:infomodel}) & $\FL\in\Mk{0}$; $\Par^{(k)}\in\Mk{k+1}$ & $\Par^{(k)}\notin\Mk{k}$ \\
consistency (Def.~\ref{def:consistent}) & $\UCB$ on $\cP$, grid $\Sq_{\Kar}$; $\Const_c$ on $\cP_c$, every grid & $\Const_c$ and $\RR$ on $\cP$ \\
\bottomrule
\end{tabular}
\end{table}

\paragraph{Numerical values} No algorithm is analysed here, so there is nothing to simulate: the
paper proves no bound that numbers could illustrate. We give instead the constants on explicit
instances, which are elementary computations from the closed forms above. For the
Bernoulli instance $\nu=(\Ber(0.7),\Ber(0.5))$, the suboptimal arm has $\Dg_a=0.2$ and, by
Proposition~\ref{prop:kinf}(iv) and~(iii),
\begin{equation*}
  2\Dg_a^2=0.0800\ \le\ \Ka_a=\kl(0.5,0.7)=0.0872\ \le\ \log\tfrac{0.5}{0.3}=0.5108,
\end{equation*}
so $1/\Ka_a=11.47$; the lower bound is nearly tight here, the upper bound is not. On the grids with
$s=2$, the first ratios $\xg_2,\dots,\xg_6$ are $1.585, 1.262, 1.161, 1.113, 1.086$ on $\Sq_2$,
$2, 1.5, 1.333, 1.25, 1.2$ on $\Ge_2$ (where $\xg_\ell=\ell/(\ell-1)$), and $2$ on $\Db_2$, whose
ends $2,4,16,256,65536,\dots$ grow doubly exponentially.

\paragraph{An instance satisfying every hypothesis at once} The instance $\nu$ above, the grid
$\Db_2$ and the learner $\Const_{c_1}$ (arm $c_1$ has law $\Ber(0.7)$) satisfy all the definitions
together: $\nu\in\cP_{c_1}\subseteq\cP$, $\Db_2$ is a grid, and $\Const_{c_1}$ belongs to $\Mk{0}$
(Proposition~\ref{prop:witnesses}(d)) and is \consistent{} of $\Db_2$ on $\cP_{c_1}$. With
$\UCB$ instead of $\Const_{c_1}$ and the grid $\Sq_2$ (here $\Kar=2$), the learner is batched on
the grid and consistent at its batch ends on the whole class $\cP$, which contains $\nu$
(Corollary~\ref{cor:ucb}).

% ======================================================================
\section{Conclusion and open questions}\label{sec:open}

\paragraph{Conclusion} We fixed the definitions of stochastic bandits on an imposed grid, with rewards in $[0,1]$: grids,
batched learners, the information models $\Mk{k}$, consistency at the batch ends, and the divergence
$\Ki$; we proved their basic relations and gave a witness and an exclusion for each. The limitations
are those of the definitions: the grid is deterministic and fixed in advance, the rewards are in
$[0,1]$, and the dual formula of $\Ki$, its continuity and the deviation inequalities for its
empirical version, which upper bounds need, are not proved here.

\paragraph{Open questions} None. The questions that motivate this setting (the price of an imposed grid for
the problem-dependent constant, in each information model) are the subject of other papers.

% ======================================================================
\appendix

\section{Useful lemmas}\label{sec:cited}

In this section we present the two results from other papers that we use.

\citet{GMS19} write $\mathbb P_1,\mathbb P_2$ for two probability measures on a measurable space
$(\Gamma,\mathcal G)$ and $Z$ for a $\mathcal G$-measurable function with values in $[0,1]$ (printed
there as $Z:\Omega\to[0,1]$, where $\Gamma$ is meant); we keep these names, except the
$\sigma$-algebra, which we write $\mathcal S$ ($\mathcal G$ is the set of grids here), with $\KL$
and $\kl$ of Definition~\ref{def:kinf}. It is used to compare the means of two laws on $[0,1]$.

% Verbatim: literature/notes/GMS19.md, Lemma 1 of Section 2.1.
\begin{lemma}[{From \citet{GMS19}, Lemma~1}]\label{cite:GMS19-lem1}
\uses{def:kinf}
Let $\mathbb P_1$ and $\mathbb P_2$ be probability measures on a measurable space
$(\Gamma,\mathcal S)$ and let $Z:\Gamma\to[0,1]$ be $\mathcal S$-measurable, with expectations
$\E_1[Z]$ and $\E_2[Z]$ under $\mathbb P_1$ and $\mathbb P_2$. Then
\begin{equation*}
  \KL(\mathbb P_1,\mathbb P_2)\ \ge\ \kl\big(\E_1[Z],\E_2[Z]\big).
\end{equation*}
\end{lemma}

\citet{ACF02} call the arms machines, write $P_i$ for the law $\nu_c$ of an arm, $\mu_i$, $\mu^*$
and $\Delta_i$ for $\mean_c$, $\ms$ and $\Dg_c$, $n$ for the number of rounds $T$, and $T_j(n)$ for
the count $N'_c(T)$ below; their regret after $n$ plays is $\mu^*n-\sum_j\mu_j\E[T_j(n)]$, which is
$\sum_c\Dg_c\E[N'_c(T)]$. In their model successive pulls of an arm give independent rewards with
its law, independently across arms; we write it with a stack of rewards per arm, as in
\citet[Section~4.6]{LS}, who state that this reward-stack model gives the same law of the history
as the canonical model and leave the proof as their Exercise~4.4. In the bound, $\pi$ is the number
$3.14159\ldots$, not a learner. UCB1 is defined there up to the order of the first $\Kar$ pulls and the
rule for ties, and the theorem holds for each choice, in particular for $\UCB$ of
Definition~\ref{def:exlearners}. It is used for Corollary~\ref{cor:ucb}.

% Verbatim: literature/notes/ACF02.md, Theorem 1 and Figure 1.
\begin{theorem}[{From \citet{ACF02}, Theorem~1}]\label{cite:ACF02-thm1}
\uses{def:instance, def:exlearners}
Let $\Kar\ge2$ and let $\nu$ be an instance with $\nu_c\in\Mu$ for every arm $c$. On a probability
space, let $(X_{c,j})_{c\in\Arms,\,j\ge1}$ be independent random variables with $X_{c,j}$ of law
$\nu_c$, and run $\UCB$ with rewards taken from these stacks: the reward of the $j$-th pull of the
arm $c$ is $X_{c,j}$. Let $N'_c(T)$ be the number of pulls of $c$ in the rounds $1,\dots,T$. Then,
for every $T\ge1$,
\begin{equation*}
  \sum_{c\in\Arms}\Dg_c\,\E\big[N'_c(T)\big]\ \le\ 8\sum_{c:\ \Dg_c>0}\frac{\log T}{\Dg_c}
  +\Big(1+\frac{\pi^2}{3}\Big)\sum_{c\in\Arms}\Dg_c .
\end{equation*}
\end{theorem}

% ======================================================================
\section{Proofs of the lemmas of Section~\ref{sec:setting}}\label{sec:pf-setting}

\subsection{Proof of Lemma~\ref{lem:counts}}\label{sec:pf-lem:counts}

\begin{proof}
\uses{def:instance, def:learner, def:counts}
\Setting $\pi$, $\nu$ with $\nu_c\in\Mu$ for all $c$, and $t\ge0$ are fixed.
\Goal (i) and (ii).

(i) Each term of $\Ng_c(t)=\sum_{s\le t}\1(\Ac_s=c)$ is in $\{0,1\}$, so $0\le\Ng_c(t)\le t$, and
going from $t$ to $t+1$ adds a nonnegative term. Since each $\Ac_s$ is exactly one arm,
$\sum_{c}\1(\Ac_s=c)=1$, and summing over $s\le t$ gives $\sum_c\Ng_c(t)=t$.

(ii) By linearity of the expectation and $\Dg_{\Ac_s}=\sum_c\Dg_c\1(\Ac_s=c)$,
\begin{align*}
  \Reg_{\nu,\pi}(t)
  &=\Enu_{\nu,\pi}\Big[\sum_{c}\Dg_c\sum_{s\le t}\1(\Ac_s=c)\Big]
   =\Enu_{\nu,\pi}\Big[\sum_{s\le t}\Dg_{\Ac_s}\Big]\\
  &=t\,\ms-\Enu_{\nu,\pi}\Big[\sum_{s\le t}\mean_{\Ac_s}\Big] && \why{$\Dg_c=\ms-\mean_c$}.
\end{align*}
Each $\Dg_{\Ac_s}$ is in $[0,\max_c\Dg_c]$ and $\max_c\Dg_c\le\ms\le1$ since the means are in
$[0,1]$, so $0\le\Reg_{\nu,\pi}(t)\le t\max_c\Dg_c\le t$.
\end{proof}

\subsection{Proof of Lemma~\ref{lem:grid}}\label{sec:pf-lem:grid}

\begin{proof}
\uses{def:grid, def:logratio}
\Setting a grid $g$ is fixed. \Goal (i)--(iii).

(i) The $\tl_\ell$ are integers and strictly increasing, so $\tl_\ell\ge\tl_{\ell-1}+1$; by induction
from $\ell=1$, $\tl_\ell\ge\tl_1+\ell-1$, which tends to infinity.

(ii) Let $t\ge1$. By (i) the set $\{\ell\ge1:\ t\le\tl_\ell\}$ is not empty, so $\bt(t)$ is defined,
and $t\le\tl_{\bt(t)}$. If $\bt(t)=1$, then $\tl_{\bt(t)-1}=\tl_0=0<t$; otherwise
$\bt(t)-1\ge1$ is not in the set by minimality, so $\tl_{\bt(t)-1}<t$. The batches are disjoint
because the intervals $(\tl_{\ell-1},\tl_\ell]$ are disjoint for a strictly increasing sequence, so
$\bt(t)$ is the only batch containing $t$. For $t=\tl_\ell$: $\tl_\ell\le\tl_\ell$ and
$\tl_{\ell-1}<\tl_\ell$, so $\bt(\tl_\ell)=\ell$. If $t\le t'$, every $\ell$ with $t'\le\tl_\ell$
satisfies $t\le\tl_\ell$, so $\bt(t)\le\bt(t')$. Finally $\bt(t)=1$ if and only if $t\le\tl_1$, by
definition of the minimum.

(iii) For $\ell\ge1$, $\tl_\ell\ge\tl_1\ge2$, so $\Lg_\ell\ge\log2>0$; $\log$ is strictly increasing,
so $\Lg_{\ell-1}<\Lg_\ell$ for $\ell\ge2$ and $\xg_\ell=\Lg_\ell/\Lg_{\ell-1}>1$. The product
telescopes:
\begin{equation*}
  \Lg_1\prod_{k=2}^{\ell}\xg_k=\Lg_1\prod_{k=2}^{\ell}\frac{\Lg_k}{\Lg_{k-1}}=\Lg_\ell .
\end{equation*}
Finally $\tl_{\ell-1}^{\,\xg_\ell}=\exp(\xg_\ell\Lg_{\ell-1})=\exp(\Lg_\ell)=\tl_\ell$.
\end{proof}

\subsection{Proof of Lemma~\ref{lem:mask}}\label{sec:pf-lem:mask}

\begin{proof}
\uses{def:batched, def:learner, def:grid}
\Setting $t\ge1$, an integer $r\ge0$ and a learner $\pi$ are fixed. \Goal measurability of
$\mask_r$ and the equivalence.

Each coordinate of $\mask_r$ is either a coordinate of $h$ or the constant $0$, so $\mask_r$ is
measurable on $\Hs_{t-1}$ with its product $\sigma$-algebra.

($\Rightarrow$) Let $h\in\Hs_{t-1}$. The histories $h$ and $\mask_r(h)$ have the same actions and the
same rewards at the rounds $s\le\min(r,t-1)$, since $\mask_r$ keeps them. If $\pi_t$ \sees{} $r$, then
$\pi_t(\cdot\mid h)=\pi_t(\cdot\mid\mask_r(h))$.

($\Leftarrow$) Let $h,h'\in\Hs_{t-1}$ with the same actions and the same rewards at the rounds
$s\le\min(r,t-1)$. Then $\mask_r(h)=\mask_r(h')$, because the other rewards are replaced by $0$ in
both. Hence
\begin{equation*}
  \pi_t(\cdot\mid h)=\pi_t(\cdot\mid\mask_r(h))=\pi_t(\cdot\mid\mask_r(h'))=\pi_t(\cdot\mid h').
\end{equation*}
The last sentence of the lemma is the case $r=\tl_{\bt(t)-1}$, for every $t$.
\end{proof}

\subsection{Proof of Lemma~\ref{lem:consist-regret}}\label{sec:pf-lem:consist-regret}

\begin{proof}
\uses{def:consistent, def:counts, def:instance, def:grid, def:learner}
\Setting $E$, $g$, $\pi$ are fixed; fix $\nu\in E$ and $\alpha>0$, and write
$u_c(\ell):=\Enu_{\nu,\pi}[\Ng_c(\tl_\ell)]/\tl_\ell^{\,\alpha}\ge0$.
\Goal $u_c(\ell)\to0$ for every suboptimal $c$ if and only if
$\Reg_{\nu,\pi}(\tl_\ell)/\tl_\ell^{\,\alpha}\to0$.

By Definition~\ref{def:counts}, and since the optimal arms have $\Dg_c=0$,
\begin{equation*}
  \frac{\Reg_{\nu,\pi}(\tl_\ell)}{\tl_\ell^{\,\alpha}}=\sum_{c:\ \Dg_c>0}\Dg_c\,u_c(\ell).
\end{equation*}
($\Rightarrow$) This is a finite sum of sequences that tend to $0$. ($\Leftarrow$) For a suboptimal
$c$, $0\le u_c(\ell)\le\Reg_{\nu,\pi}(\tl_\ell)/(\Dg_c\,\tl_\ell^{\,\alpha})$, which tends to $0$.
\end{proof}

% ======================================================================
\section{Proof of Theorem~\ref{thm:main}}\label{sec:pf-thm:main}

\begin{proof}[Proof of Theorem~\ref{thm:main}]
\proves{thm:main}
\uses{def:instance, def:learner, def:counts, def:grid, def:exgrids, def:batched, def:infomodel, def:consistent, def:seqcons, def:exlearners, lem:grid, lem:consist-regret, prop:witnesses}
\Setting the set of arms $\Arms$, with $\Kar\ge2$. \Goal items (i)--(iii) of
Theorem~\ref{thm:main}.

\emph{Idea.} Every item reduces to comparing the last batch end before a round in two grids.
Agreeing on more ends implies agreeing on fewer, which gives the inclusions of (i); strictness is
the parity learner of Proposition~\ref{prop:witnesses}(d). If $g$ is finer than $g'$, the last end
of $g$ before a round is at least that of $g'$, which gives (ii), and the ends of $g'$ are a
subsequence of those of $g$. On $\Sq_s$ the last end before a round $t>s$ is $t-1$, so batching
only constrains the first batch, and the ends form a tail of all rounds, which gives (iii) with
Lemma~\ref{lem:consist-regret}. The proof follows the items: (i) inclusions, strictness and
components; (ii) batching and consistency; (iii) batching and consistency.

\textbf{Item (i), inclusions.} Let $k\le k'$ and $\Pi\in\Mk{k}$. Condition (a) of
Definition~\ref{def:infomodel} does not depend on $k$. For (b) at level $k'$: if two grids $g,g'$
agree on $\tl_j$ for all $j\le\ell+k'$, they agree for all $j\le\ell+k$, since $\ell+k\le\ell+k'$ (with
$\ell+\infty=\infty$); so (b) at level $k$ gives $\pi^g_t=\pi^{g'}_t$. Hence $\Pi\in\Mk{k'}$.

\textbf{Item (i), strictness.} Let $k<k'$; then $k\in\Nat$ and $k+1\le k'$. By
Proposition~\ref{prop:witnesses}(d), $\Par^{(k)}\in\Mk{k+1}\subseteq\Mk{k'}$ and
$\Par^{(k)}\notin\Mk{k}$, so $\Mk{k}\ne\Mk{k'}$.

\textbf{Item (i), components.} If $\Pi\in\Mk{\infty}$, then $\pi^g$ is batched on $g$ by (a). Conversely,
let $\pi$ be batched on $g$, and let $\Pi$ have $\pi^g:=\pi$ and $\pi^{g'}:=\RR$ for $g'\ne g$. Each
component is batched on its grid ($\RR$ ignores the rewards, so it sees them up to any $r\ge0$), so (a)
holds, and (b) is void for $k=\infty$ because its hypothesis forces $g=g'$. Hence $\Pi\in\Mk{\infty}$.

\textbf{Item (ii), batching.} Let $g=(\tl_j)_j$ be finer than $g'=(\tl'_j)_j$, let $\pi$ be batched on
$g'$, and let $t\ge1$. Write $j:=\bt_g(t)-1$ and $i:=\bt_{g'}(t)-1$ (Definition~\ref{def:grid}: the
batch of $t$ in $g$ and in $g'$, minus one) for the indices of the last ends
before $t$ (Lemma~\ref{lem:grid}(ii): $\tl_j<t$ and $\tl'_i<t$). The end $\tl'_i$ of $g'$ is an end
of $g$ (also when $i=0$, as $\tl'_0=\tl_0=0$), and it is smaller than $t$; since $\tl_j$ is the
largest end of $g$ smaller than $t$ (the next one, $\tl_{j+1}$, is $\ge t$),
\begin{equation*}
  \tl'_i\le\tl_j .
\end{equation*}
Now let $h,h'$ have the same actions and the same rewards at the rounds $s\le\min(\tl_j,t-1)$. They
have the same rewards at the rounds $s\le\min(\tl'_i,t-1)$, so $\pi_t(\cdot\mid h)=\pi_t(\cdot\mid h')$
because $\pi_t$ sees the rewards up to $\tl'_i$. Hence $\pi_t$ sees the rewards up to $\tl_j$, for
every $t$: $\pi$ is batched on $g$.

\textbf{Item (ii), consistency.} The ends of $g'$ form a subset of the ends of $g$; since both
sequences are strictly increasing, there is a strictly increasing $\varphi$ with
$\tl'_\ell=\tl_{\varphi(\ell)}$ for $\ell\ge1$. For $\nu\in E$, a suboptimal $c$ and $\alpha>0$, the
sequence $\ell\mapsto\Enu_{\nu,\pi}[\Ng_c(\tl'_\ell)]/(\tl'_\ell)^\alpha$ is a subsequence of
$\ell\mapsto\Enu_{\nu,\pi}[\Ng_c(\tl_\ell)]/\tl_\ell^{\,\alpha}$, which tends to $0$; so does the
subsequence.

\textbf{Item (iii), batching.} On $\Sq_s$, $\tl_\ell=s+\ell-1$. For $t\le s=\tl_1$,
$\bt(t)=1$ (Lemma~\ref{lem:grid}(ii)) and $\tl_{\bt(t)-1}=0$: the condition of
Definition~\ref{def:batched} at $t$ is that $\pi_t$ ignores the rewards. For $t>s$,
\begin{equation*}
  \tl_{t-s}=t-1<t\le t=\tl_{t-s+1},
\end{equation*}
so $\bt(t)=t-s+1$ and $\tl_{\bt(t)-1}=t-1$; a policy $\pi_t$ always sees the rewards up to $t-1$,
since its histories have length $t-1$. Hence $\pi$ is batched on $\Sq_s$ if and only if $\pi_t$
ignores the rewards for every $t\le s$.

\textbf{Item (iii), consistency.} Fix $\nu\in E$, $\alpha>0$ and write
$r(T):=\Reg_{\nu,\pi}(T)/T^\alpha$. By Lemma~\ref{lem:consist-regret}, $\pi$ is consistent at the
batch ends of $\Sq_s$ on $E$ if and only if $r(s+\ell-1)\to0$ as $\ell\to\infty$ for all such $\nu$
and $\alpha$. The map $\ell\mapsto s+\ell-1$ is a bijection from $\{\ell\ge1\}$ onto $\{T\ge s\}$,
and the limit of a sequence depends only on its tail, so $r(s+\ell-1)\to0$ if and only if $r(T)\to0$
as $T\to\infty$, which is sequential consistency (Definition~\ref{def:seqcons}).
\end{proof}

% ======================================================================
\section{Proof of Proposition~\ref{prop:kinf}}\label{sec:pf-prop:kinf}

We first record two properties of $\kl$, used for the lower bound in (iii) and for (iv).

\begin{lemma}[monotonicity and Pinsker for $\kl$]\label{lem:kl}
\uses{def:kinf}
Let $p\in[0,1)$. Then $q\mapsto\kl(p,q)$ is nondecreasing on $[p,1]$ and continuous on $[p,1)$,
with $\kl(p,p)=0$, and
\begin{equation*}
  \kl(p,q)\ge2(q-p)^2\qquad\text{for every }q\in[p,1].
\end{equation*}
\end{lemma}

\begin{proof}
\uses{def:kinf}
\Setting $p\in[0,1)$. \Goal monotonicity, continuity, $\kl(p,p)=0$ and the quadratic bound.

At $q=1$, $\kl(p,1)=+\infty$ since $p<1$, which is larger than every other value. On $[p,1)$, if $p=0$
then $\kl(0,q)=\log\frac1{1-q}$, which is continuous, nondecreasing and $0$ at $q=0$, and
differentiable on $(0,1)$ with
$\frac{\partial}{\partial q}\kl(0,q)=\frac1{1-q}=\frac{q-0}{q(1-q)}$. If $p>0$,
then $[p,1)\subseteq(0,1)$, where $\kl(p,\cdot)$ is finite and differentiable, $\kl(p,p)=0$, and
\begin{equation*}
  \frac{\partial}{\partial q}\kl(p,q)=-\frac pq+\frac{1-p}{1-q}=\frac{q-p}{q(1-q)}\ \ge0\qquad(q\in[p,1)).
\end{equation*}
So in both cases $\frac{\partial}{\partial q}\kl(p,q)=\frac{q-p}{q(1-q)}$ on $(p,1)$.
For the quadratic bound, let $f(q):=\kl(p,q)-2(q-p)^2$ on $[p,1)$: it is continuous, $f(p)=0$, and on
$(p,1)$, since $q(1-q)\le1/4$,
\begin{equation*}
  f'(q)=(q-p)\Big(\frac1{q(1-q)}-4\Big)\ \ge0 .
\end{equation*}
Hence $f\ge0$ on $[p,1)$, and the bound holds at $q=1$ because $\kl(p,1)=+\infty$.
\end{proof}

\begin{proof}[Proof of Proposition~\ref{prop:kinf}]
\proves{prop:kinf}
\uses{def:instance, def:kinf, cite:GMS19-lem1, lem:kl}
\Setting $P\in\Mu$, $p=\mean(P)\in[0,1]$, $\mu\in\R$. \Goal (i)--(v).

\emph{Idea.} Items (i) and (ii) read off the definition: no law on $[0,1]$ has mean above $1$, and
$P$ itself is admissible when $\mu<p$. For (iii), every admissible $Q$ has mean $q>\mu$, and
data processing to the reward gives $\KL(P,Q)\ge\kl(p,q)\ge\kl(p,\mu)$; Pinsker's inequality
(Lemma~\ref{lem:kl}) gives the quadratic bound. The upper bound moves just enough mass of $P$ to
the point $1$: the mixture $(1-u)P+u\delta_1$ has a density bounded by $1/(1-u)$. Item (iv) uses the
Bernoulli laws as competitors, and (v) is (iii) at $\mu=\ms$.

\textbf{(i)} Every $Q\in\Mu$ has $\mean(Q)\le1\le \mu$, so the set in the definition of $\Ki(P,\mu)$ is
empty and $\Ki(P,\mu)=+\infty$.

\textbf{(ii)} If $\mu<p$, then $Q:=P\in\Mu$ has $\mean(Q)>\mu$ and $\KL(P,P)=0$, so $\Ki(P,\mu)=0$.

\textbf{(iii), lower bound.} Let $p\le \mu<1$ and let $Q\in\Mu$ with $q:=\mean(Q)>\mu$. Apply
Lemma~\ref{cite:GMS19-lem1} on $(\Gamma,\mathcal S):=(\R,\text{Borel})$ with $\mathbb P_1:=P$,
$\mathbb P_2:=Q$ and $Z(x):=\min(\max(x,0),1)$, which is measurable with values in $[0,1]$. Since
$P$ and $Q$ give mass $1$ to $[0,1]$, where $Z(x)=x$, we have $\E_1[Z]=p$ and $\E_2[Z]=q$, so
\begin{align*}
  \KL(P,Q)&\ge\kl(p,q) && \why{Lemma~\ref{cite:GMS19-lem1}}\\
          &\ge\kl(p,\mu) && \why{Lemma~\ref{lem:kl}: $\kl(p,\cdot)$ nondecreasing on $[p,1]$, $p\le \mu<q\le1$}.
\end{align*}
Taking the infimum over such $Q$ gives $\Ki(P,\mu)\ge\kl(p,\mu)$, and $\kl(p,\mu)\ge2(\mu-p)^2$ by
Lemma~\ref{lem:kl} (here $p\le \mu<1$, so $p<1$).

\textbf{(iii), upper bound.} Let $u_0:=\frac{\mu-p}{1-p}\in[0,1)$ (well defined since $p\le \mu<1$) and
let $u\in(u_0,1)$. The measure $Q_u:=(1-u)P+u\,\delta_1$ is in $\Mu$, with mean
\begin{equation*}
  \mean(Q_u)=(1-u)p+u=p+u(1-p)>p+u_0(1-p)=\mu .
\end{equation*}
Since $Q_u\ge(1-u)P$, $P$ is absolutely continuous with respect to $Q_u$ and its density satisfies
$\frac{dP}{dQ_u}\le\frac1{1-u}$, $Q_u$-almost everywhere, hence $P$-almost everywhere. Therefore
\begin{equation*}
  \Ki(P,\mu)\le\KL(P,Q_u)=\int\log\frac{dP}{dQ_u}\,dP\le\log\frac1{1-u}.
\end{equation*}
Letting $u$ decrease to $u_0$ gives $\Ki(P,\mu)\le\log\frac1{1-u_0}=\log\frac{1-p}{1-\mu}$, which is
finite.

\textbf{(iv)} Let $P=\Ber(p)$ with $p\le \mu<1$. The lower bound $\Ki(P,\mu)\ge\kl(p,\mu)$ is (iii). For
the upper bound, every $q\in(\mu,1)$ gives $\Ber(q)\in\Mu$ with mean $q>\mu$, so
$\Ki(P,\mu)\le\KL(\Ber(p),\Ber(q))=\kl(p,q)$. As $q$ decreases to $\mu$, $\kl(p,q)\to\kl(p,\mu)$ by the
continuity of Lemma~\ref{lem:kl} on $[p,1)$, so $\Ki(P,\mu)\le\kl(p,\mu)$.

\textbf{(v)} Let $\nu\in\cP$ and $a$ suboptimal. Apply (iii) with $P:=\nu_a\in\Mu$, $p=\mean_a$ and
$\mu:=\ms$: indeed $\mean_a<\ms<1$. It gives
$2\Dg_a^2=2(\ms-\mean_a)^2\le \Ka_a\le\log\frac{1-\mean_a}{1-\ms}<\infty$, and $2\Dg_a^2>0$.
\end{proof}

% ======================================================================
\section{Proofs of Proposition~\ref{prop:witnesses} and Corollary~\ref{cor:ucb}}\label{sec:pf-prop:witnesses}

\subsection{Proof of Proposition~\ref{prop:witnesses}}

\begin{proof}[Proof of Proposition~\ref{prop:witnesses}]
\proves{prop:witnesses}
\uses{def:instance, def:grid, def:logratio, def:exgrids, def:learner, def:counts, def:batched, def:infomodel, def:consistent, def:exlearners, lem:grid, lem:counts}
\Setting $\Kar\ge2$ and $s\ge\max(2,\Kar)$. \Goal items (a)--(e).

\emph{Idea.} Each item is a direct check on an explicit object. A learner whose arm depends on the
round only ignores the rewards, hence is batched on every grid; the frozen leader reads only the
rewards up to the last end, and two histories with opposite rewards make it pull different arms.
For the information models, two grids that agree on the ends a learner reads give it the same
policy, and the parity learner is separated by two grids that differ by one at the end
$\tl_{k+1}$. For consistency, a learner that never pulls a suboptimal arm has zero counts, and a
learner that pulls a fixed suboptimal arm a constant fraction of the time has linear counts.

\textbf{(a)} $\Ber(p_c)$ gives mass $1$ to $\{0,1\}\subseteq[0,1]$ and has mean $p_c$, so
$\ms=\max_cp_c<1$ and the instance is in $\cP$. An arm $\delta_1$ has mean $1$, so $\ms=1$; an arm
$\mathcal N(0,1)$ gives mass $<1$ to $[0,1]$, so it is not in $\Mu$. Neither instance is in $\cP$.

\textbf{(b)} In the three cases $\tl_0=0$, $\tl_1=s\ge\max(2,\Kar)$ and the sequence is strictly
increasing for $\ell\ge1$: $s+\ell-1$ increases by $1$, $s2^{\ell-1}$ doubles, and
$s^{2^{\ell}}=(s^{2^{\ell-1}})^2>s^{2^{\ell-1}}$ since $s^{2^{\ell-1}}\ge2$. The ratios are, for $\ell\ge2$,
\begin{equation*}
  \Sq_s:\ \xg_\ell=\frac{\log(s+\ell-1)}{\log(s+\ell-2)},\qquad
  \Ge_s:\ \xg_\ell=\frac{\log s+(\ell-1)\log2}{\log s+(\ell-2)\log2},\qquad
  \Db_s:\ \xg_\ell=\frac{2^{\ell-1}\log s}{2^{\ell-2}\log s}=2 .
\end{equation*}
For $\Ge_s$ the numerator and the denominator differ by $\log2$ and tend to infinity, so
$\xg_\ell\to1$. For $\Sq_s$, $\log(s+\ell-1)-\log(s+\ell-2)=\log(1+\frac1{s+\ell-2})\to0$ while the
denominator tends to infinity, so $\xg_\ell\to1$. The sequence $\tl_\ell:=\ell$ has $\tl_1=1<2$, so
it is not a grid.

\textbf{(c)} $\RR$ and $\Const_c$ ignore the rewards (their arm at round $t$ depends on $t$ only), so
they see the rewards up to every $r\ge0$, and are batched on every grid.

Let $g$ be a grid and $t\ge1$. If $t\le\tl_1$, $\FL^g$ plays as $\RR$, which ignores the rewards. If
$t>\tl_1$, then $\ell=\bt(t)-1\ge1$ and $\tl_\ell<t$ (Lemma~\ref{lem:grid}(ii)); the arm pulled is
a function of the actions and of the rewards $y_s$, $s\le\tl_\ell$, so $\FL^g_t$ sees the rewards up
to $\tl_\ell=\tl_{\bt(t)-1}$. Hence $\FL^g$ is batched on $g$.

$\FL^g$ uses the rewards at round $t:=\tl_1+1$: take the history $h$ of the first $\tl_1$ rounds of
$\RR$ (each arm is pulled at least once there, as $\tl_1\ge\Kar$) with $y_s:=\1(a_s=c_1)$, and the
history $h'$ with the same actions and $y'_s:=\1(a_s=c_2)$. Then $\FL^g$ pulls $c_1$ after $h$ and
$c_2$ after $h'$, so $\FL^g_t$ does not ignore the rewards.

$\React$ is batched on no grid: for every grid, $2\le\tl_1$, so $\bt(2)=1$ and $\React_2$ must ignore
the rewards. The histories $h:=(c_1,1)$ and $h':=(c_1,0)$ have the same action, and $\React$ pulls $c_1$
after $h$ and $c_2\ne c_1$ after $h'$.

\textbf{(d), $\FL\in\Mk{0}$.} Condition (a) of Definition~\ref{def:infomodel} is (c). For (b) with
$k=0$, let $g,g'$ agree on $\tl_1,\dots,\tl_\ell$ and let $t$ be in batch $\ell+1$ of both. Then
$\bt_g(t)=\bt_{g'}(t)=\ell+1$. If $\ell=0$, both policies are those of $\RR$ at $t$. If $\ell\ge1$,
both pull the first maximiser of the frozen empirical means computed from the rounds
$s\le\tl_\ell=\tl'_\ell$, the same function of the history. Hence $\FL^g_t=\FL^{g'}_t$.

\textbf{(d), $(\Const_c)_{g\in\Gr}\in\Mk{0}$.} Condition (a) is (c): $\Const_c$ is batched on every
grid. Condition (b) holds for every $k$, because the policy $\Const_c$ at a round $t$ is the same for
all grids (the Dirac mass at $c$ after every history).

\textbf{(d), $\Par^{(k)}\in\Mk{k+1}$.} Each $\Par^{(k),g}$ ignores the rewards, so it is batched on
$g$. Let $g,g'$ agree on $\tl_j$ for $1\le j\le\ell+k+1$, and let $t$ be in batch $\ell+1$ of both.
Then $\bt_g(t)=\bt_{g'}(t)=\ell+1$ and $\tl_{\ell+1+k}=\tl'_{\ell+1+k}$, so the two policies pull the
same arm at $t$.

\textbf{(d), $\Par^{(k)}\notin\Mk{k}$.} Let $g:=\Sq_s$ ($\tl_j=s+j-1$ for $j\ge1$)
and $g'$ the sequence with $\tl'_0:=0$, $\tl'_j:=s+j-1$ for $1\le j\le k$ and $\tl'_j:=s+j$ for
$j\ge k+1$. Then $g'$ is a grid: it is strictly increasing (it increases by $2$ from $j=k$ to $j=k+1$
when $k\ge1$, and by $1$ elsewhere) and $\tl'_1\ge s\ge\max(2,\Kar)$. Take $\ell:=0$ and $t:=1$:
round $1$ is in batch $1$ of both grids, since $\tl_1,\tl'_1\ge2$, and the grids agree on $\tl_j$
for $1\le j\le0+k$. But
\begin{equation*}
  \tl_{\bt(1)+k}=\tl_{k+1}=s+k,\qquad \tl'_{\bt(1)+k}=\tl'_{k+1}=s+k+1
\end{equation*}
have different parities, so $\Par^{(k),g}$ pulls one of $c_1,c_2$ at round $1$ and $\Par^{(k),g'}$
pulls the other. Condition (b) of $\Mk{k}$ fails.

\textbf{(e)} Let $\nu\in\cP_c$. Every suboptimal arm $c'$ of $\nu$ differs from $c$ because
$\mean_c=\ms$, so under $\Const_c$, $\Ng_{c'}(t)=0$ for all $t$: $\Const_c$ is consistent at the batch
ends of every grid on $\cP_c$.

On $\cP$: fix a grid. Let $\nu$ have $\nu_c:=\Ber(1/4)$ and $\nu_{c'}:=\Ber(1/2)$ for $c'\ne c$;
then $\nu\in\cP$ by (a) and $c$ is suboptimal. Under $\Const_c$, $\Ng_c(\tl_\ell)=\tl_\ell$, so with
$\alpha:=1/2$,
\begin{equation*}
  \frac{\Enu_{\nu,\Const_c}[\Ng_c(\tl_\ell)]}{\tl_\ell^{1/2}}=\tl_\ell^{1/2}\to\infty
\end{equation*}
by Lemma~\ref{lem:grid}(i). Under $\RR$, every block of $\Kar$ consecutive rounds pulls each arm
once, so $\Ng_c(t)\ge\lfloor t/\Kar\rfloor\ge t/\Kar-1$ for all $t$ and every arm $c$ (with $\Ng_c(t)\le t$,
Lemma~\ref{lem:counts}(i)); with the same $\nu$ and $\alpha$,
\begin{equation*}
  \frac{\Enu_{\nu,\RR}[\Ng_c(\tl_\ell)]}{\tl_\ell^{1/2}}\ge\frac{\tl_\ell}{\Kar\,\tl_\ell^{1/2}}-\frac1{\tl_\ell^{1/2}}\to\infty .
\end{equation*}
Neither learner is consistent at the batch ends on $\cP$.
\end{proof}

\subsection{Proof of Corollary~\ref{cor:ucb}}\label{sec:pf-cor:ucb}

\begin{proof}[Proof of Corollary~\ref{cor:ucb}]
\proves{cor:ucb}
\uses{def:instance, def:learner, def:counts, def:seqcons, def:exgrids, def:batched, def:consistent, def:exlearners, cite:ACF02-thm1, thm:main}
\Setting $\Kar\ge2$, so that $s:=\Kar$ satisfies $s\ge\max(2,\Kar)$ and $\Sq_{\Kar}$ is the grid
$\tl_\ell=\Kar+\ell-1$. \Goal $\UCB$ is batched on $\Sq_{\Kar}$, and consistent at its batch ends on
$\cP$.

\emph{Idea.} By Theorem~\ref{thm:main}(iii) with $s=\Kar$, both properties reduce to sequential
ones: the first $\Kar$ policies must ignore the rewards, which $\UCB$ does, and $\UCB$ must be
sequentially consistent on $\cP$, which follows from the logarithmic regret bound of
Theorem~\ref{cite:ACF02-thm1} once the stack model of that theorem is shown to have the law
$\Pnu_{\nu,\UCB}$ of Definition~\ref{def:learner}. Steps: (1) the law; (2) sequential consistency;
(3) conclusion.

\textbf{Step 1: the stack model gives the law $\Pnu_{\nu,\UCB}$.} This identification is standard:
\citet[Section~4.6]{LS} state it, and its proof is their Exercise~4.4. We write it out for a
deterministic learner. Fix $\nu$ with $\nu_c\in\Mu$ for all $c$. $\UCB$ is deterministic: its arm at round $t$ is $f_t(h)$ for a measurable map $f_t$ on
$\Hs_{t-1}$. In the stack model of Theorem~\ref{cite:ACF02-thm1}, let $A'_t:=f_t(H'_{t-1})$,
$R'_t:=X_{A'_t,\,N'_{A'_t}(t-1)+1}$ and $H'_t:=(A'_1,R'_1,\dots,A'_t,R'_t)$. We show by induction on
$t$ that $H'_t$ has the law of $H_t$ under $\Pnu_{\nu,\UCB}$; for $t=0$ both are the point of
$\Hs_0$. Let $t\ge1$, an arm $c$ and an integer $j\ge1$, and let $X^{(c,j)}$ be the family of
stacks in which $X_{c,i}$ is replaced by $0$ for every $i\ge j$. Run $\UCB$ on $X^{(c,j)}$, with
history $H^{(c,j)}$ and counts $N^{(c,j)}$. Up to round $t-1$, the two runs read the same rewards
as long as $c$ has been pulled at most $j-1$ times; so, by induction on the rounds, the events
$\{N'_c(t-1)=j-1\}$ and $\{N^{(c,j)}_c(t-1)=j-1\}$ coincide and $H'_{t-1}=H^{(c,j)}_{t-1}$ on them.
These are functions of $X^{(c,j)}$, which is independent of $X_{c,j}$. Hence, for every bounded
measurable $\varphi$ on $\Hs_{t}$,
\begin{equation*}
  \E\big[\varphi(H'_{t-1},c,X_{c,j})\,\1\{A'_t=c,\ N'_c(t-1)=j-1\}\big]
  =\E\Big[\int\varphi(H'_{t-1},c,y)\,\nu_c(dy)\,\1\{A'_t=c,\ N'_c(t-1)=j-1\}\Big],
\end{equation*}
since $A'_t=f_t(H'_{t-1})$. On the event of the indicator, $R'_t=X_{c,j}$. Summing over $c$ and
$j\ge1$,
\begin{equation*}
  \E\big[\varphi(H'_t)\big]=\E\Big[\int\varphi\big(H'_{t-1},f_t(H'_{t-1}),y\big)\,\nu_{f_t(H'_{t-1})}(dy)\Big].
\end{equation*}
The same identity holds for $H_t$ under $\Pnu_{\nu,\UCB}$, by the two conditional laws of
Definition~\ref{def:learner}. With the induction hypothesis, $H'_t$ and $H_t$ have the same law. The
count of an arm up to round $T$ is a function of the history, so
$\E[N'_c(T)]=\Enu_{\nu,\UCB}[\Ng_c(T)]$ for all $c$ and $T$.

\textbf{Step 2: $\UCB$ is sequentially consistent on $\cP$.} Let $\nu\in\cP$ and $\alpha>0$. By
Step~1, Definition~\ref{def:counts} and Theorem~\ref{cite:ACF02-thm1}, for every $T\ge1$,
\begin{equation*}
  0\le\frac{\Reg_{\nu,\UCB}(T)}{T^\alpha}\le
  \frac{8\log T}{T^\alpha}\sum_{c:\ \Dg_c>0}\frac1{\Dg_c}+\frac1{T^\alpha}\Big(1+\frac{\pi^2}{3}\Big)\sum_{c}\Dg_c
  \ \xrightarrow[T\to\infty]{}\ 0,
\end{equation*}
since $\log T/T^\alpha\to0$ (here $\pi=3.14159\ldots$ is the number). This is Definition~\ref{def:seqcons} with $E=\cP$.

\textbf{Step 3: conclusion.} For $t\le\Kar$, $\UCB$ pulls $c_t$ whatever the history, so its policy
at $t$ ignores the rewards. By Theorem~\ref{thm:main}(iii) with $s=\Kar$, $\UCB$ is batched on
$\Sq_{\Kar}$; by the same item with $E=\cP$ and Step~2, it is consistent at the batch ends of
$\Sq_{\Kar}$ on $\cP$.
\end{proof}

% ======================================================================
\section{Sources of the techniques}\label{sec:attribution}

\begin{center}
\begin{tabular}{@{}p{0.45\textwidth}p{0.45\textwidth}@{}}
\toprule
Technique (where) & Earliest source found \\
\midrule
$\mathcal K_{\inf}$ as an infimum of $\mathrm{KL}$ over a model, $\inf\emptyset=+\infty$ (Definition~\ref{def:kinf}) & \citet{BK96}; in this form \citet[Section~1.2]{GMS19} \\
Consistency as $o(T^\alpha)$ regret, or counts (Definitions~\ref{def:seqcons}, \ref{def:consistent}) & \citet{LR85}; \citet[Definition~16.1]{LS}; \citet[Definition~1]{GMS19} \\
$\KL\ge\kl$ of the means for laws on $[0,1]$, hence $\Ki\ge\kl$ (Proposition~\ref{prop:kinf}(iii)) & \citet[Section~6.1]{CGMMS13}, from Hoeffding's bound on the moment-generating function; in the form used here (data processing with a $[0,1]$-valued $Z$), \citet[Lemma~1]{GMS19} \\
The chain $\Ki\ge\kl\ge2(\cdot)^2$ (Proposition~\ref{prop:kinf}(iii)) & \citet[(7)]{GHMS22} \\
Mixing with an atom, $(1-u)P+u\delta_1$ (Proposition~\ref{prop:kinf}(iii)) & \citet[Appendix~B.1]{GHMS22} (mixing with $\delta_x$) \\
Upper bound $\Ki(P,\mu)\le\log\frac{1-p}{1-\mu}$ (Proposition~\ref{prop:kinf}(iii)) & this form: source not known; the uniform bound $-\log(1-\mu)$: \citet[Lemma~14]{HT10} \\
Pinsker's inequality for Bernoulli laws with the constant $2$, and its proof by integrating $\partial_q\kl$ (Lemma~\ref{lem:kl}) & Csisz\'ar (1967), Kullback (1967) and Kemperman (1969), and Sch\"utzenberger (1953) in the binary case, after the survey of \citet[Sections~3--4]{Rio23}; the proof is Csisz\'ar's, as written in \citet[(16)]{Rio23} \\
Learners that know the whole imposed grid or only its past ends (Definition~\ref{def:infomodel}, $k=\infty$ and $k=0$) & \citet{QW26} (batch-aware and batch-oblivious learners) \\
Intermediate information models $\Mk{k}$, $1\le k<\infty$ (Definition~\ref{def:infomodel}) & source not known; anytime learners on a static grid: \citet{JTXHXG21} \\
Batched learners as policies that see the rewards up to the last end (Definition~\ref{def:batched}) & \citet{PRCS16} for chosen grids; \citet{QW26} for imposed ones \\
A consistent learner on the sequential grid (Corollary~\ref{cor:ucb}) & \citet[Theorem~1]{ACF02} \\
The reward-stack model has the law of the canonical interaction (Corollary~\ref{cor:ucb}, Step~1) & stated in \citet[Section~4.6]{LS}, proof left as their Exercise~4.4 (hint: Doob's optional stopping theorem); the argument written here replaces a stack by zeros from an index on instead \\
\bottomrule
\end{tabular}
\end{center}

% ======================================================================
\section{Notation}\label{sec:notation}

Every symbol is a link to its definition.
\begin{center}
\begin{tabular}{@{}p{0.27\textwidth}p{0.44\textwidth}p{0.23\textwidth}@{}}
\toprule
Symbol & Meaning & Defined in \\
\midrule
$\Arms$, $\Kar$ & set of arms, number of arms & Definition~\ref{def:instance} \\
$\Mu$, $\mean(Q)$ & laws on $[0,1]$, mean of $Q$ & Definition~\ref{def:instance} \\
$\mean_c$, $\ms$, $\Dg_c$ & mean of arm $c$, largest mean, gap & Definition~\ref{def:instance} \\
$\cP$, $\cP_c$ & instances with rewards in $[0,1]$ and $\ms<1$; those where $c$ is optimal & Definitions~\ref{def:instance}, \ref{def:exlearners} \\
$\KL$, $\Ber$, $\kl$ & relative entropy, Bernoulli law, Bernoulli relative entropy & Definition~\ref{def:kinf} \\
$\Ki(P,\mu)$, $\Ka_a$ & infimum of $\KL$ over $\Mu$ above the level $\mu$; $\Ki(\nu_a,\ms)$ & Definition~\ref{def:kinf} \\
$\Hs_t$, $\pi_t$ & histories of length $t$, policy at round $t$ & Definition~\ref{def:learner} \\
$\Ac_t$, $\Rw_t$, $\Pnu_{\nu,\pi}$, $\Enu_{\nu,\pi}$ & arm, reward, law and expectation of the interaction & Definition~\ref{def:learner} \\
$\Ng_c(t)$, $\Reg_{\nu,\pi}(t)$ & pulls of $c$ up to round $t$, regret & Definition~\ref{def:counts} \\
$\Gr$, $\tl_\ell$, $\Bg_\ell$, $\bt(t)$, $\bt_g(t)$ & grids, batch ends, batch sizes, batch of round $t$ (in the grid $g$) & Definition~\ref{def:grid} \\
$\Lg_\ell$, $\xg_\ell$ & $\log\tl_\ell$, ratio $\Lg_\ell/\Lg_{\ell-1}$ & Definition~\ref{def:logratio} \\
$\Sq_s$, $\Ge_s$, $\Db_s$ & sequential, geometric, doubly exponential grids & Definition~\ref{def:exgrids} \\
$\mask_r$ & history with the rewards after round $r$ set to $0$ & Definition~\ref{def:batched} \\
$\Mk{k}$ & information model & Definition~\ref{def:infomodel} \\
$\RR$, $\Const_c$, $\React$, $\FL$, $\Par^{(k)}$, $\UCB$ & example learners & Definition~\ref{def:exlearners} \\
\bottomrule
\end{tabular}
\end{center}

\paragraph{Vocabulary} A policy \sees{} $r$, or \ignores{} (Definition~\ref{def:batched}); a learner
is \batched{} on a grid (Definition~\ref{def:batched}); a grid is \finer{} than another
(Definition~\ref{def:grid}); a learner is \consistent{} (Definition~\ref{def:consistent}) or
\seqcons{} (Definition~\ref{def:seqcons}).

\bibliographystyle{plainnat}
\bibliography{refs,../../papers/library}

\end{document}
